%! Piecewise \& Step Functions — Unit 1, Lesson 1.4 — Exit Ticket
\documentclass[10pt]{article}
\usepackage{algebra2-article}
\usepackage{algebra2-boxes}
\newcommand{\TallMath}[1]{$\displaystyle #1\rule[-1.4em]{0pt}{3.2em}$}
% Coordinate grid: {xmin}{xmax}{ymin}{ymax}
\newcommand{\lgrid}[4]{%
  \draw[step=1, linegray!40, very thin] (#1,#3) grid (#2,#4);
  \draw[->, semithick, charcoal] (#1,0)--(#2,0) node[below right, font=\scriptsize]{$x$};
  \draw[->, semithick, charcoal] (0,#3)--(0,#4) node[left, font=\scriptsize]{$y$};}
% Endpoint dots
\newcommand{\closeddot}[2]{\fill[forest] (#1,#2) circle (3.5pt);}
\newcommand{\opendot}[2]{\draw[very thick, forest, fill=white] (#1,#2) circle (3.5pt);}

\begin{document}

\pageheader{Unit 1, Lesson 1.4}{Exit Ticket}

\vspace{0.10in}

No notes. Show your work in the space provided.

\begin{enumerate}[leftmargin=1.8em, itemsep=14pt]
  \item \textbf{Evaluate.} For $f(x)=\begin{cases} 3x, & x<1 \\ x+4, & x\ge 1 \end{cases}$
        \\[4pt]
        $f(0)=\blank{1.3cm}$ \qquad $f(1)=\blank{1.3cm}$ \qquad $f(3)=\blank{1.3cm}$
        \\[4pt]
        Which piece did you use for $f(1)$, and how do you know? \writelines{2}

  \item \textbf{Read the graph.} The piecewise function below is shown.
        \begin{center}
        \begin{tikzpicture}[scale=0.5]
          \lgrid{-3}{3}{-2}{4}
          \draw[very thick, forest] (-3,-1)--(0,-1);
          \opendot{0}{-1}
          \draw[very thick, forest] (0,1)--(2,3);
          \closeddot{0}{1}
        \end{tikzpicture}
        \end{center}
        $f(-2)=\blank{1.0cm}$; \quad $f(0)=\blank{1.0cm}$. \quad Is $f$ continuous at $x=0$? Explain.
        \writelines{2}

  \item \textbf{Multiple choice (SOL-style).} The greatest-integer function gives $\lfloor -1.5\rfloor=$
        \begin{enumerate}[label=(\Alph*), leftmargin=1.8em, itemsep=2pt]
          \item $-1$
          \item $-2$
          \item $1$
          \item $2$
        \end{enumerate}
\end{enumerate}

\end{document}
